\documentclass[11pt,a4paper]{article}
\usepackage[T1]{fontenc}
\usepackage[utf8]{inputenc}
\usepackage{lmodern}
\usepackage[margin=19mm,top=17mm,bottom=18mm]{geometry}
\usepackage{amsmath,amssymb,bm}
\usepackage{microtype}
\usepackage{hyperref}
\usepackage{titlesec}
\hypersetup{hidelinks,pdftitle={Learning local structure and motion from simulation},pdfauthor={}}
\titleformat{\section}{\large\bfseries}{}{0pt}{}
\titlespacing*{\section}{0pt}{10pt}{5pt}
\titlespacing*{\paragraph}{0pt}{7pt}{0.6em}
\setlength{\parindent}{0pt}
\setlength{\parskip}{4pt}
\setlength{\abovedisplayskip}{6pt}
\setlength{\belowdisplayskip}{6pt}
\setlength{\emergencystretch}{1em}
\clubpenalty=10000
\widowpenalty=10000
\newcommand{\R}{\mathbb R}
\newcommand{\norm}[1]{\left\lVert #1\right\rVert}
\newcommand{\avg}[1]{\left\langle #1\right\rangle}
\newcommand{\x}{\mathbf x}
\newcommand{\X}{\mathbf X}
\newcommand{\Pp}{\mathbf P}
\newcommand{\z}{\mathbf z}
\newcommand{\Z}{\mathbf Z}
\newcommand{\cstate}{\mathbf c}
\newcommand{\Y}{\mathbf Y}
\newcommand{\Sim}{\operatorname{Sim}}

\begin{document}
\begin{center}
{\LARGE\bfseries Learning local structure and motion from simulation}
\end{center}
\vspace{-4pt}

We formulate a local model of structure and motion with an evolving memory.
Observed positions and velocities initialize a representation of a central
atom and its neighbourhood, together with a summary of the observed past.
The representation and memory evolve jointly and yield predictions of local
atomic states and structural measurements. Molecular dynamics under a fixed
interatomic potential supplies future trajectories and their responses to
small initial-state perturbations for joint training.

\section{Starting snapshots and physical time}

\paragraph{Indices.}
The dataset contains $N$ starting snapshots. Snapshot $i$ contains $n_i$ atoms
and is observed at physical time $t_i$. We use
\[
\begin{aligned}
i&=1,\ldots,N &&\text{starting snapshot},\\
b&&&\text{MD continuation from that snapshot},\\
j&=1,\ldots,n_i &&\text{central atom}.
\end{aligned}
\]
The index $i$ identifies a dataset example; its associated time is $t_i$.
Different examples may come from different source MD trajectories.

\paragraph{Saved observations and prediction times.}
Let $\delta t$ be the numerical integration step of MD. If observations are
saved every $m$ integration steps, their spacing is $\Delta t=m\delta t$.
Within one uniformly saved source trajectory, let $t_{\mathrm{start}}$ be its
first saved time and $\operatorname{frame}(i)$ the integer index of snapshot $i$.
Then
\[
t_i=t_{\mathrm{start}}+\operatorname{frame}(i)\,\Delta t.
\]
A prediction starts at $t_i$. Let $h$ be the time advanced by one learned
step, $K$ the number of steps, and $\tau_{\max}=Kh$ the total prediction horizon.
We use $\tau\in[0,\tau_{\max}]$ for elapsed time after the starting snapshot.
Step $k$ corresponds to elapsed time $\tau_k$ and physical time $s_{ik}$:
\[
\boxed{\tau_k=kh,\qquad s_{ik}=t_i+\tau_k,\qquad k=0,\ldots,K.}
\]
The physical continuation saves observations at the times $s_{ik}$.
For example, $\delta t=1\,\mathrm{fs}$, $h=5\,\mathrm{fs}$ and $K=20$
give a $0.1\,\mathrm{ps}$ prediction, with preceding history sampled at
$\Delta t=0.1\,\mathrm{ps}$.

\paragraph{Complete starting state.}
At physical time $s$, atom $\ell$ has position $\mathbf r_{i\ell}(s)$ and
velocity $\mathbf v_{i\ell}(s)$. Its mass is $m_\ell$, and its momentum is
$\mathbf p_{i\ell}(s)=m_\ell\mathbf v_{i\ell}(s)$. The complete positions
and momenta are
\[
\X_i(s)=\bigl(\mathbf r_{i1}(s),\ldots,\mathbf r_{in_i}(s)\bigr),
\qquad
\Pp_i(s)=\bigl(\mathbf p_{i1}(s),\ldots,\mathbf p_{in_i}(s)\bigr).
\]
The pair $(\X_i(t_i),\Pp_i(t_i))$, together with the periodic cell, initializes
the physical continuation.

\section{Local state and observed history}

\paragraph{Tracked neighbourhood.}
Let $\mathcal N_{ij}$ be the atoms selected around centre $j$ at time $t_i$,
including the centre. Their identities and ordering are retained throughout
the observed history and predicted future.
For each $\ell\in\mathcal N_{ij}$, define the relative position and velocity
in the nonperiodic case by
\[
\mathbf q_{ij\ell}(s)=\mathbf r_{i\ell}(s)-\mathbf r_{ij}(s),
\qquad
\mathbf w_{ij\ell}(s)=\mathbf v_{i\ell}(s)-\mathbf v_{ij}(s).
\]
When a displacement crosses the simulation box boundary, we account for the
periodic boundary conditions.
Collect these vectors in the arrays
$\x_{ij}(s)=(\mathbf q_{ij\ell}(s))_{\ell\in\mathcal N_{ij}}$ and
$\mathbf w_{ij}(s)=(\mathbf w_{ij\ell}(s))_{\ell\in\mathcal N_{ij}}$.
The local atomic state is
\[
\boxed{S_{ij}(s)=\bigl(\x_{ij}(s),\mathbf w_{ij}(s)\bigr).}
\]

\paragraph{Observed history.}
Let $L$ be the number of available observations, including the present.
We write $C_{ij}$ for the preceding history:
\[
C_{ij}=\bigl(S_{ij}(t_i-(L-1)\Delta t),\ldots,S_{ij}(t_i-\Delta t)\bigr).
\]
The present state is written separately as $S_{ij,0}=S_{ij}(t_i)$.
For four observations, the model therefore receives the states at
$t_i-3\Delta t$, $t_i-2\Delta t$, $t_i-\Delta t$ and $t_i$.
The input consists of these ordered local positions and velocities.
The times specify which observations are supplied and which future is predicted.

\section{Physical continuations and measurements}

\paragraph{Hamiltonian dynamics.}
Let $U(\X)$ be the energy of the complete configuration under the fixed
interatomic potential. With momenta $\Pp=(\mathbf p_1,\ldots,\mathbf p_{n_i})$,
the Hamiltonian is
\[
\boxed{H(\X,\Pp)=U(\X)+\sum_{\ell=1}^{n_i}
\frac{\norm{\mathbf p_\ell}^2}{2m_\ell}.}
\]
The short teacher trajectory evolves in the starting periodic cell, held fixed
throughout the continuation. Its equations of motion are
\[
\frac{d\mathbf r_\ell}{ds}=\frac{\partial H}{\partial\mathbf p_\ell}
=\frac{\mathbf p_\ell}{m_\ell},
\qquad
\frac{d\mathbf p_\ell}{ds}=-\frac{\partial H}{\partial\mathbf r_\ell}.
\]
The simulator integrates these equations for the complete system:
\[
\Gamma_{ib}(t_i:t_i+\tau_{\max})
=\Sim\!\left(\X_i(t_i),\Pp_i(t_i);\tau_{\max}\right).
\]
With fixed initial positions, momenta and cell, this continuation is
deterministic. The teacher supplies one such continuation per starting state.
Interactions with surrounding atoms contribute to the evolution of the local
neighbourhood.

\paragraph{Local future measurements.}
From continuation $b$, extract the local state at each prediction time:
\[
S_{ibj,k}=S_{ibj}(s_{ik})
=\bigl(\x_{ibj}(s_{ik}),\mathbf w_{ibj}(s_{ik})\bigr).
\]
Let $\phi_{\mathrm{str}}$ compute structural measurements from local positions,
and let $\phi_{\mathrm{vel}}$ compute motion measurements from local positions
and relative velocities. Their concatenation defines a fixed measurement map
$\phi$ with $M$ outputs:
\[
\boxed{
\Y_{ibj,k}=\phi(S_{ibj,k})
=\begin{pmatrix}
\phi_{\mathrm{str}}(\x_{ibj}(s_{ik}))\\
\phi_{\mathrm{vel}}(\x_{ibj}(s_{ik}),\mathbf w_{ibj}(s_{ik}))
\end{pmatrix}\in\R^M.
}
\]
Structural measurements include bond order, topological descriptors and smooth
local density descriptors. Motion measurements include relative-velocity
moments, local flow and instantaneous rates of change of bond order.
The sequence $(\Y_{ibj,0},\ldots,\Y_{ibj,K})$ describes how these properties
change along one physical continuation.

\section{Local dynamics and memory}

\paragraph{Eliminating the surrounding degrees of freedom.}
The local state $S$ retains a subset of the complete atomic configuration.
Atoms outside the neighbourhood influence its motion and are themselves
affected by it. Eliminating these degrees of freedom produces dependence on
the local past. Choose a probability distribution over complete initial
states and use conditional averaging over states with the same local
observation. For fixed $i,b,j$, write $S(\tau)=S_{ibj}(t_i+\tau)$.
Let $A(S)$ be the conditional average of the instantaneous local rate of
change, $K_{\mathrm{mem}}$ the memory function, and $\eta_{ibj}(\tau)$ the
contribution from the unresolved initial state. The Mori--Zwanzig identity
then has the form~\cite{lin2021,gupta2025}
\[
\boxed{
\frac{dS(\tau)}{d\tau}
=A(S(\tau))
+\int_0^\tau K_{\mathrm{mem}}\bigl(\tau-\tau',S(\tau')\bigr)\,d\tau'
+\eta_{ibj}(\tau).
}
\]
Here $\tau'$ is an earlier elapsed time. The memory function describes
delayed coupling through the omitted degrees of freedom. It may depend
nonlinearly on the past local state. The unresolved contribution depends on
the complete initial condition; its statistical description generally has
temporal correlations. The averaging distribution and projection determine
all three terms.

\paragraph{Information in the observed past.}
The history $C_{ij}$ provides information about the surroundings through
their previous effects on the neighbourhood. A finite observed history
constrains the hidden starting conditions. We approximate the predictable
evolution with a learned memory state initialized from this history.
The deterministic model below leaves the remaining distribution of
unresolved forcing unspecified. Its value losses fit local trajectories
under the distribution of physical training states.

\section{Proposed learned evolution with memory}

\paragraph{Initial representation.}
Let $E=(E_z,E_c)$ be an encoder with two outputs. The vector $\z$ contains
features of the tracked atoms, and $\cstate$ contains $d_c$ memory components.
The combined learned state is $\Z=(\z,\cstate)$. Its initial value is
\[
\boxed{
\Z_{ij,0}=E(S_{ij,0},C_{ij})
=\bigl(E_z(S_{ij,0},C_{ij}),E_c(S_{ij,0},C_{ij})\bigr).
}
\]
Attention combines current atomic features with the ordered observations in
$C_{ij}$. The observed history remains fixed during prediction; the memory
$\cstate$ evolves with the predicted trajectory.

\paragraph{Continuous evolution.}
Let $f(\z,\cstate)$ be a learned rate of change of $\z$, and let $g_r(\z)$
write information into memory component $r$, for $r=1,\ldots,d_c$.
Let $\lambda_r>0$ be its learned decay rate, with units of inverse time.
We use the coupled equations
\[
\boxed{
\frac{d\z_{ij}}{d\tau}=f(\z_{ij},\cstate_{ij}),
\qquad
\frac{dc_{ij,r}}{d\tau}=-\lambda_r c_{ij,r}+g_r(\z_{ij}).
}
\]
The second equation implies
\[
c_{ij,r}(\tau)
=e^{-\lambda_r\tau}c_{ij,r}(0)
+\int_0^\tau e^{-\lambda_r(\tau-\tau')}
g_r(\z_{ij}(\tau'))\,d\tau'.
\]
Each component therefore summarizes earlier states over a characteristic
duration $1/\lambda_r$. This finite-memory construction follows the
auxiliary-state approach used in iLED~\cite{menier2025}. The functions
$f$ and $g$ provide learned nonlinear coupling between the state and its
memory. Decay rates are shared across vector components that rotate together.

Let $G_h$ denote numerical integration of the coupled equations over a
duration $h$. At the prediction times,
\[
\boxed{
\Z_{ij,k}=(\z_{ij}(\tau_k),\cstate_{ij}(\tau_k)),
\qquad
\Z_{ij,k+1}=G_h(\Z_{ij,k}).
}
\]
The evolution is autonomous: $f$ and $g$ receive the learned state.
Physical timestamps organize the observations and evaluation times.

\paragraph{Atomic and measurement decoders.}
Let $D$ decode a latent state into local positions and velocities, using the
observed history. Predictions are anchored at the actual initial state:
\[
\boxed{
\widehat S_{ij,k}=S_{ij,0}
+D(\z_{ij,k},C_{ij})-D(\z_{ij,0},C_{ij}).
}
\]
Thus $\widehat S_{ij,0}=S_{ij,0}$. Write
$\widehat S_{ij,k}=(\widehat\x_{ij,k},\widehat{\mathbf w}_{ij,k})$
for its position and velocity components.

Let $F$ be a second decoder that predicts the measurement vector directly
from the latent state:
\[
\widehat\Y_{ij,k}=F(\z_{ij,k}).
\]
Its present output is fitted to the measured present state, and its future
outputs are fitted to the simulated measurement sequence. The measurement
targets supervise the learned representation. Rotating the atomic input
rotates vector features and predicted atomic vectors; the scalar measurement
outputs remain invariant.

\paragraph{Trajectory and measurement losses.}
For a vector of errors $\mathbf e$ and a positive diagonal weight matrix $W$,
define $\norm{\mathbf e}_W^2=\mathbf e^{\mathsf T}W\mathbf e$.
Matrices $W_S$ and $W_Y$ set the component scales for atomic states and
measurements, using training statistics and fixed group weights. Angle brackets
denote an average over training examples and the sampled time indices, with
their specified weights. The value losses are
\[
\begin{aligned}
L_{\mathrm{state}}&=
\avg{\norm{\widehat S_{ij,k}-S_{ibj,k}}_{W_S}^2}_{k\geq1},\\
L_{\mathrm{meas}}&=
\avg{\norm{\widehat\Y_{ij,k}-\Y_{ibj,k}}_{W_Y}^2}_{k\geq0}.
\end{aligned}
\]
To train the changes along the same connected trajectory, compare consecutive
measurement increments:
\[
L_{\mathrm{inc}}=
\avg{\norm{
(\widehat\Y_{ij,k+1}-\widehat\Y_{ij,k})
-(\Y_{ibj,k+1}-\Y_{ibj,k})}_{W_Y}^2}.
\]
Each pair of consecutive targets belongs to the same starting state,
continuation and central atom.

\paragraph{State consistency.}
Reconstruction fits $D(E_z(S_{ibj,k},C_{ij}),C_{ij})$ to
$S_{ibj,k}$.
Kinematic consistency fits the predicted displacement between steps to the
integrated relative velocity:
\[
\widehat\x_{ij,k+1}-\widehat\x_{ij,k}
\simeq\frac h2\bigl(\widehat{\mathbf w}_{ij,k+1}
+\widehat{\mathbf w}_{ij,k}\bigr).
\]
Let $L_{\mathrm{cons}}$ collect the weighted squared reconstruction and
kinematic errors. Latent-state consistency will also contribute to this loss.

\section{Learning the response to initial perturbations}

\paragraph{Perturbed physical states.}
Choose position and momentum directions $\mathbf d^X_{ij}$ and
$\mathbf d^P_{ij}$ for the complete system, with zero entries outside the
selected neighbourhood. Their entries have units of position and momentum.
For a dimensionless amplitude $a$, the perturbed initial state is
\[
\X_i^{(j,a)}(t_i)=\X_i(t_i)+a\mathbf d^X_{ij},
\qquad
\Pp_i^{(j,a)}(t_i)=\Pp_i(t_i)+a\mathbf d^P_{ij}.
\]
A position query sets $\mathbf d^P_{ij}=0$; a velocity query changes momenta
according to the chosen velocity direction and atomic masses.
The simulator propagates these initial changes through the complete system.
Let $S_{ibj,k}^{(a)}$ be the resulting local state. Its response is
\[
\mathbf u_{ibj,k}
=\left.\frac{\partial S_{ibj,k}^{(a)}}{\partial a}\right|_{a=0}.
\]
This vector gives the change in local positions and velocities per unit
perturbation amplitude at time $s_{ik}$. Differentiating the numerical
integration steps propagates this response to every saved point along the
short trajectory. Atom identities and local periodic-image choices are fixed
for the derivative calculation.

\paragraph{Response of measurements.}
Let $\phi_{\mathrm{sm}}$ contain the differentiable density and velocity
measurements. For a differentiable function $f$, its Jacobian $J_f$ maps a
small input change to the corresponding first-order output change.
The physical response of these measurements is
\[
\boxed{
\mathbf R_{ibj,k}
=\left.\frac{\partial\phi_{\mathrm{sm}}(S_{ibj,k}^{(a)})}{\partial a}\right|_{a=0}
=J_{\phi_{\mathrm{sm}}}(S_{ibj,k})\,\mathbf u_{ibj,k}.
}
\]
Bond-order and topological measurements enter through the value losses.

\paragraph{Response of the learned trajectory.}
Apply the same initial local perturbation to the neural prediction, keeping
the observed past $C_{ij}$ fixed. Both outputs of $E$ are evaluated on the
perturbed present state, so the perturbation may change the initial memory.
Let $\Z_{ij,k}^{(a)}=(\z_{ij,k}^{(a)},\cstate_{ij,k}^{(a)})$ and
$\widehat S_{ij,k}^{(a)}$ be the resulting latent and decoded trajectories,
including the perturbed initial state in the anchoring formula.
Let $F_{\mathrm{sm}}$ select the density and velocity outputs of $F$. The
predicted responses are
\[
\widehat{\mathbf u}_{ij,k}
=\left.\frac{\partial\widehat S_{ij,k}^{(a)}}{\partial a}\right|_{a=0},
\qquad
\widehat{\mathbf R}_{ij,k}
=\left.\frac{\partial F_{\mathrm{sm}}(\z_{ij,k}^{(a)})}{\partial a}\right|_{a=0}.
\]
Let $W_u$ and $W_R$ scale the atomic and measurement response errors using
response magnitudes measured on training trajectories. The response losses are
\[
L_{\mathrm{atomic\,resp}}
=\avg{\norm{\widehat{\mathbf u}_{ij,k}-\mathbf u_{ibj,k}}_{W_u}^2},
\qquad
L_{\mathrm{measured\,resp}}
=\avg{\norm{\widehat{\mathbf R}_{ij,k}-\mathbf R_{ibj,k}}_{W_R}^2}.
\]
The averages include sampled perturbation directions and future steps. These
losses match the direction and magnitude of the physical response throughout
the predicted trajectory. Differentiation passes through the coupled state
and memory dynamics at every learned step.

\paragraph{Transport of response in latent space.}
The physical continuation defines a teacher trajectory in the learned
coordinates. For elapsed time $\tau$, define
\[
\z^{\mathrm{enc}}_{ibj}(\tau)
=E_z(S_{ibj}(t_i+\tau),C_{ij}).
\]
Its memory $\cstate^{\mathrm{enc}}_{ibj}(\tau)$ is driven by this entire
encoded physical trajectory:
\[
\frac{dc^{\mathrm{enc}}_{ibj,r}}{d\tau}
=-\lambda_r c^{\mathrm{enc}}_{ibj,r}
+g_r(\z^{\mathrm{enc}}_{ibj}),
\qquad
\cstate^{\mathrm{enc}}_{ibj}(0)=E_c(S_{ij,0},C_{ij}).
\]
Thus $\Z^{\mathrm{enc}}_{ibj,k}
=(\z^{\mathrm{enc}}_{ibj}(\tau_k),\cstate^{\mathrm{enc}}_{ibj}(\tau_k))$
contains an encoded physical state and a memory accumulated along its
continuation. This construction also applies to each perturbed teacher.
Write $\mathbf e^z_{ibj}(\tau)$ and $\mathbf e^c_{ibj}(\tau)$ for the
derivatives of its state and memory with respect to the initial perturbation
amplitude. Their sampled values form the combined tangent
\[
\mathbf e_{ibj,k}
=\left.\frac{\partial\Z^{\mathrm{enc},(a)}_{ibj,k}}
{\partial a}\right|_{a=0}
=\bigl(\mathbf e^z_{ibj,k},\mathbf e^c_{ibj,k}\bigr).
\]
The state component is
$\mathbf e^z_{ibj,k}=J_{E_z}(S_{ibj,k};C_{ij})\mathbf u_{ibj,k}$.
The Jacobian differentiates the current-state argument with the history
fixed. The memory component propagates the responses of all preceding
teacher states:
\[
\frac{de^c_{ibj,r}}{d\tau}
=-\lambda_r e^c_{ibj,r}
+J_{g_r}(\z^{\mathrm{enc}}_{ibj})\mathbf e^z_{ibj},
\qquad
\mathbf e^c_{ibj}(0)=J_{E_c}(S_{ij,0};C_{ij})\mathbf u_{ibj,0}.
\]
These tangent functions follow the same physical continuation. Learned
parameters are held fixed when differentiating with respect to $a$.

Let $J_{G_h}$ differentiate the learned step with respect to its combined
input state. Consistency of values and responses requires
\[
\boxed{\begin{aligned}
G_h(\Z^{\mathrm{enc}}_{ibj,k})
&\simeq\Z^{\mathrm{enc}}_{ibj,k+1},\\
J_{G_h}(\Z^{\mathrm{enc}}_{ibj,k})\,\mathbf e_{ibj,k}
&\simeq\mathbf e_{ibj,k+1}.
\end{aligned}}
\]
The first residual contributes to $L_{\mathrm{cons}}$, using training-based
scales for the state and memory components. The second compares a response
advanced by the learned dynamics with the response of the encoded physical
continuation, including its accumulated memory.
Let $W_T$ scale latent response errors by the mean squared magnitude of the
encoded physical response. The transport loss is
\[
L_{\mathrm{transport}}
=\avg{\norm{
J_{G_h}(\Z^{\mathrm{enc}}_{ibj,k})\mathbf e_{ibj,k}
-\mathbf e_{ibj,k+1}}_{W_T}^2}.
\]
The scale is held fixed during each gradient calculation. Sampled transitions
apply this constraint across the physical continuation.

\newpage
\section{Joint training and use of the representation}

Let $\lambda_Y$, $\lambda_{\mathrm{inc}}$, $\lambda_u$, $\lambda_R$,
$\lambda_T$ and $\lambda_C$ be nonnegative weights that set the contributions
of the measurement, increment, response, transport and consistency losses.
The joint training objective is
\[
\boxed{\begin{aligned}
L={}&L_{\mathrm{state}}
+\lambda_Y L_{\mathrm{meas}}
+\lambda_{\mathrm{inc}}L_{\mathrm{inc}}
+\lambda_C L_{\mathrm{cons}}\\
&+\lambda_u L_{\mathrm{atomic\,resp}}
+\lambda_R L_{\mathrm{measured\,resp}}
+\lambda_T L_{\mathrm{transport}}.
\end{aligned}}
\]
The parameters of $E$, $f$, $g$, $D$, $F$ and the positive decay rates
$\lambda_r$ are optimized together. The potential and measurement functions
remain fixed. The physical simulator generates training targets; prediction
advances the learned state and memory and applies the decoders. Derivative
supervision constrains the response over the simulated teacher horizon.
Learning longer memory effects requires physical observations over the
corresponding durations.

After training, the learned model can be frozen. A readout is a separate
predictive model fitted to features of its initial combined state or predicted
state and memory trajectory. Rotation-invariant summaries provide local features for
tasks such as nucleus formation. Comparing these readouts with readouts of
the directly measured input history assesses the predictive information
provided by the learned dynamics. Source trajectories define the training
and test groups, and normalization is fitted on training data. Training
budgets and checkpoint rules are fixed before test evaluation.

Trajectory evaluation includes atomic and measurement errors, connected
increments and temporal correlations. Response errors are compared with the
error of predicting zero response. A matched model without an explicit memory
state provides a control for the contribution of $\cstate$. State-dependent
memory models for molecular dynamics~\cite{ge2024} and memory-augmented neural
operators~\cite{buitrago2025} motivate testing these temporal properties.

\begin{thebibliography}{9}
\bibitem{lin2021}
Y.~T. Lin, Y. Tian, M. Anghel, and D. Livescu.
\emph{Data-driven learning for the Mori--Zwanzig formalism:
A generalization of the Koopman learning framework}.
SIAM Journal on Applied Dynamical Systems, 2021.
\href{https://doi.org/10.1137/21M1401759}{doi:10.1137/21M1401759}.

\bibitem{gupta2025}
P. Gupta, P.~J. Schmid, D. Sipp, T. Sayadi, and G. Rigas.
\emph{Mori--Zwanzig latent space Koopman closure for nonlinear autoencoder}.
Proceedings of the Royal Society A, 2025.
\href{https://doi.org/10.1098/rspa.2024.0259}{doi:10.1098/rspa.2024.0259}.

\bibitem{menier2025}
E. Menier, S. Kaltenbach, M. Yagoubi, M. Schoenauer, and P. Koumoutsakos.
\emph{Interpretable learning of effective dynamics for multiscale systems}.
Proceedings of the Royal Society A, 2025.
\href{https://doi.org/10.1098/rspa.2024.0167}{doi:10.1098/rspa.2024.0167}.

\bibitem{ge2024}
P. Ge, Z. Zhang, and H. Lei.
\emph{Data-driven learning of the generalized Langevin equation with
state-dependent memory}.
Physical Review Letters, 133:077301, 2024.
\href{https://doi.org/10.1103/PhysRevLett.133.077301}{doi:10.1103/PhysRevLett.133.077301}.

\bibitem{buitrago2025}
R. Buitrago Ruiz, T. Marwah, A. Gu, and A. Risteski.
\emph{On the benefits of memory for modeling time-dependent PDEs}.
International Conference on Learning Representations, 2025.
\href{https://openreview.net/forum?id=o9kqa5K3tB}{OpenReview}.
\end{thebibliography}

\end{document}
